\section*{Bab \ref{ch:b2:phanerozoic-earth} --- Bumi Fanerozoikum}

\begin{solution}{exo:b2:phanerozoic-earth:1}
(a) Paris (\ang{5.1}~LU) dan New York (\ang{3.6}~LS); Cape Town dan Sydney terletak dekat \ang{60}~LS.
(b) Cape Town: $61.9 - 33.9 = \ang{28.0}$, $28.0 \times 111.2 = \qty{3114}{km}$; New York:
$40.7 + 3.6 = \ang{44.3}$, \qty{4926}{km}. (c) $3114/300 = \qty{10.4}{km}$ per juta tahun,
\qty{1.0}{cm} per tahun; $4926/300 = \qty{16.4}{km}$ per juta tahun, \qty{1.6}{cm} per tahun.
\end{solution}

\begin{solution}{exo:b2:phanerozoic-earth:2}
(a) 72, 120.1, 50.5, dan \qty{135.4}{Ma}; rata-rata $(444 - 66)/4 = \qty{94.5}{Ma}$. (b) Krisis
akhir Perm (Siberia, berselisih \qty{1.9}{Ma}), krisis akhir Trias (CAMP, \qty{0.4}{Ma}), dan krisis
akhir Kapur (Deccan, \qty{0.5}{Ma}). (c) Tidak: selangnya berkisar dari 50 sampai
\qty{135}{Ma}; rata-rata selalu ada untuk himpunan tanggal apa pun.
\end{solution}

\begin{solution}{exo:b2:phanerozoic-earth:3}
(a) $\qty{91}{m}/\qty{0.036}{mm.yr^{-1}} = \num{2.5e6}$ tahun. (b) $3.7/0.036 = 103$ kali. (c)
Jauh dari pantai mana pun, di air yang jernih, hanya sedikit lumpur dari sungai yang mencapai dasar; hujan
cangkang alga adalah sedimen utamanya.
\end{solution}

\begin{solution}{exo:b2:phanerozoic-earth:4}
(a) Dari batu pasir sungai sampai kapur, endapannya makin dalam ke atas: suatu transgresi.
Pasir pantai di atas kapur merekam suatu regresi. (b) Tidak: itu merekam kenaikan muka laut relatif,
yang juga dapat dihasilkan oleh penurunan tempat itu atau berkurangnya pasokan sedimen. (c)
Menemukan transgresi yang sama, pada waktu yang sama, di benua-benua lain, atau menghilangkan
penurunan itu dengan \omterm{def:b2:basins-compared:backstripping}{pengupasan balik}.
\end{solution}

\begin{solution}{exo:b2:phanerozoic-earth:5}
(a) $233/97 = \qty{2.4}{m}$ per juta tahun, \qty{2.4}{\micro m} per tahun. (b) $3.7/0.0024 =
1540$: kenaikan masa kini sekitar 1500 kali lebih cepat. (c) Kurva yang dikupas balik: pengeboran telah menunjukkan
bahwa diagram siklus benar tentang waktu penurunan tetapi melebih-lebihkan besarnya, dengan faktor dua atau
lebih; New Jersey memiliki biasnya sendiri (topografi dinamis), sehingga perkiraannya $100 \pm
\qty{50}{m}$.
\end{solution}

\begin{solution}{exo:b2:phanerozoic-earth:6}
(a) $\num{e5} \times \num{e9} \times 2900 = \qty{2.9e17}{kg}$ basal, dan $\num{2.9e17}
\times 0.0075 = \qty{2.2e15}{kg}$ karbon dioksida, $\num{2.2e15}/0.04401 = \qty{4.9e16}{mol}$.
(b) $\num{278e-6} \times (44.01/28.97) \times \num{5.1e18} = \qty{2.15e15}{kg}$. (c) Satu denyut
melepaskan kira-kira sebanyak yang dikandung seluruh udara pada 1750. (d) $\num{e17}/\num{2.15e15} = 46.5$: hampir lima puluh.
\end{solution}

\begin{solution}{exo:b2:phanerozoic-earth:7}
(a) $2630 \times \frac{\pi}{6} \times (\num{e4})^3 = \qty{1.38e15}{kg}$. (b) $\tfrac12 \times
\num{1.38e15} \times (\num{11.2e3})^2 = \qty{8.7e22}{J}$; pada \qty{20}{km.s^{-1}},
\qty{2.76e23}{J}. (c) Satu tahun kehilangan \qty{47}{TW} adalah $\num{47e12} \times \num{3.156e7} =
\qty{1.48e21}{J}$: 59 dan 186 tahun.
\end{solution}

\begin{solution}{exo:b2:phanerozoic-earth:8}
(a) Warna gelap, laminasi halus, tanpa liang, banyak bahan organik, pirit; sering ada pergeseran
isotop karbon dan tanda-tanda pemanasan. (b) Peristiwa Toarsium dengan Karoo--Ferrar, OAE~1a dengan
Ontong Java. (c) Penanggalan membawa ketidakpastian sampai jutaan tahun, dan ada banyak peristiwa dan
provinsi, sehingga kebetulan cukup mungkin terjadi; kaitan sebab-akibat memerlukan provinsi yang mulai
tepat sebelum peristiwa itu dan suatu penjejak vulkanisme, seperti merkuri, dalam sedimen yang sama.
\end{solution}

\begin{solution}{exo:b2:phanerozoic-earth:9}
(a) $60.9 + 5.1 = \ang{66.0}$, $66.0 \times 111.2 = \qty{7339}{km}$ dalam \qty{200}{Ma}:
\qty{36.7}{km} per juta tahun, \qty{3.7}{cm} per tahun. (b) Lebih dari dua kali
\qty{1.6}{cm} per tahun sejak \qty{300}{Ma}. (c) Hanya perubahan lintang yang diketahui; perubahan
bujur, yang bebas dalam kerangka paleomagnetik, dan setiap jalan memutar menambah jaraknya.
\end{solution}

\begin{solution}{exo:b2:phanerozoic-earth:10}
(a) Fraksi dasar samudra yang berumur $t$ adalah $p(t) = \frac2T(1 - t/T)$. Maka $\int_0^T tp\,\mathrm{d}t =
\frac2T\bigl(\frac{T^2}2 - \frac{T^2}3\bigr) = \frac T3$ dan $\int_0^T \sqrt t\,p\,\mathrm{d}t =
\frac2T\bigl(\frac23 T^{3/2} - \frac25 T^{3/2}\bigr) = \frac8{15}\sqrt T$. (b) $T = 3 \times 64.2 =
\qty{192.6}{Ma}$; $2500 + \frac8{15} \times 350 \times 13.878 = \qty{5091}{m}$. (c) $T/2 =
\qty{96.3}{Ma}$: \qty{4332}{m}; penurunannya \qty{759}{m}, \qty{2}{\percent} lebih kecil daripada
\qty{774}{m}. (d) Model merata: $\frac23\sqrt\tau = \frac23\sqrt{2\bar t} = 0.943\sqrt{\bar t}$;
segitiga: $\frac8{15}\sqrt{3\bar t} = 0.924\sqrt{\bar t}$. Jawabannya jauh lebih bergantung pada umur rata-rata
daripada pada bentuk sebarannya.
\end{solution}

\begin{solution}{exo:b2:phanerozoic-earth:11}
(a) Dalam $\Delta s$, hanya faktor $1 - f^{-1/2}$ yang bergantung pada $f$; bagilah dengan nilainya pada $f = 2$,
$1 - 2^{-1/2} = 0.2929$. (b) $f = (1 - 0.2929\,\Delta s/432)^{-2}$: 1.071 untuk \qty{50}{m}, 1.239
untuk \qty{150}{m}; 7 sampai \qty{24}{\percent} lebih cepat daripada laju kini. (c) $432 \times (1 -
0.8^{-1/2})/0.2929 = \qty{-174}{m}$: laut akan turun \qty{174}{m}. (d) Sebaran umur
berubah hanya ketika dasar samudra tua tersubduksi dan digantikan, yang memakan waktu berorde $\tau$,
sekitar \qty{100}{Ma}; laut mengikutinya selama puluhan juta tahun.
\end{solution}

\begin{solution}{exo:b2:phanerozoic-earth:12}
(a) $2^{-7.4/2} = 0.077$: \qty{92}{\percent} hilang. (b) $t = 2\log_2(1/0.19) = \qty{4.8}{Ma}$.
(c) $4.8/0.060 = 80$ kali lebih cepat. (d) Dalam jenjang yang panjang, kepunahan latar saja melenyapkan
sebagian besar spesies, sehingga denyut yang singkat hanya menambah jumlah total yang sudah besar dan sukar dipisahkan darinya;
ketidaklengkapan catatan juga mengaburkan kemunculan terakhir ke arah yang lebih tua.
\end{solution}

\begin{solution}{pb:b2:phanerozoic-earth:1}
\textbf{1.} 2500, 4250, 5304, dan \qty{6000}{m}.
\textbf{2.} $\tau = 2 \times 64.2 = \qty{128.4}{Ma}$; $2500 + 350\sqrt{128.4} = \qty{6466}{m}$.
\textbf{3.} $\num{293e6}/(65000 \times 128.4) = \qty{35.1}{km}$ per juta tahun,
\qty{35.1}{mm} per tahun: laju yang lambat.
\textbf{4.} $\frac1\tau\int_0^\tau(2500 + 350\sqrt t)\,\mathrm{d}t = 2500 + \frac23 \times
350\sqrt\tau = \qty{5144}{m}$.
\textbf{5.} Hukum akar kuadrat terlalu dalam untuk dasar samudra yang lebih tua daripada beberapa puluh juta tahun,
yang mendatar; umur yang sebenarnya tidak tersebar merata, karena dasar samudra muda paling melimpah; dan
sedimen, \omterm{def:b2:oceanic-islands:seamount}{gunung laut}, serta dataran tinggi mengangkat dasar samudra.
\textbf{6.} $\num{293e6} \times 5.144 = \qty{1.51e9}{km^3}$, \qty{13}{\percent} lebih banyak daripada
seluruh samudra: model itu terlalu dalam.
\textbf{7.} $v = \qty{17.6}{km}$ per juta tahun; $17.6 \times 128.4 = \qty{2260}{km}$ di
setiap sisi.
\textbf{8.} $t = 1000/17.6 = \qty{56.8}{Ma}$; $2500 + 350\sqrt{56.8} = \qty{5138}{m}$.
\textbf{9.} Pada satu sayap, dengan $x = vt$: $\int_0^{v\tau}\bigl(d(\tau) - d(x/v)\bigr)\mathrm{d}x
= vc\int_0^\tau(\sqrt\tau - \sqrt t)\,\mathrm{d}t = \frac13 vc\,\tau^{3/2}$; dua kalinya untuk kedua sayap.
$\frac23 \times 17.6 \times 0.350 \times 128.4^{3/2} = \frac23 \times 17.6 \times 0.350 \times
1454.9 = \qty{5975}{km^2}$.
\textbf{10.} $65000 \times 5975 = \qty{3.88e8}{km^3}$.
\textbf{11.} $\num{3.88e8}/\num{1.33e9} = 0.29$.
\textbf{12.} Dasar samudranya rata-rata lebih muda, jadi lebih panas, kurang padat, dan lebih tinggi; jalur dasar samudra
muda di sepanjang sumbu lebih lebar.
\textbf{13.} $\tau = \qty{64.2}{Ma}$; $2500 + \frac23 \times 350\sqrt{64.2} = \qty{4370}{m}$.
\textbf{14.} $5144 - 4370 = \qty{774}{m}$.
\textbf{15.} $\num{293e6} \times 0.774 = \qty{2.27e8}{km^3}$.
\textbf{16.} $\num{2.27e8}/\num{362e6} = \qty{0.627}{km}$, \qty{627}{m}.
\textbf{17.} $\Delta s = 627 \times (3300 - 1025)/3300 = \qty{432}{m}$.
\textbf{18.} $f = \bigl(1 - 0.2929 \times 100/432\bigr)^{-2} = 1.15$: \qty{15}{\percent} lebih cepat.
\textbf{19.} $0.149 \times 148 = \qty{22.1}{\mega km^2}$.
\textbf{20.} $48.1 + 0.32 \times 7.7 = \qty{50.6}{\percent}$, \qty{74.9}{\mega km^2}.
\textbf{21.} Ketika naik, laut menutupi luas yang lebih besar, sehingga volume yang sama membentuk lapisan yang lebih tipis
(\cref{prop:b2:phanerozoic-earth:flooding}).
\textbf{22.} Luas yang tergenang pada 100, 200, 300, dan \qty{400}{m}: 22.1, 41.7, 57.3, dan
\qty{71.2}{\mega km^2}. Langkah-langkah: $(362 + 11.0) \times 0.1 = \qty{37.3e6}{km^3}$, lalu 39.4, 41.2,
dan \qty{42.6e6}{km^3}; kumulatifnya 37.3, 76.7, 117.9, dan \qty{160.5e6}{km^3}. Volume yang harus ditempatkan
adalah \qty{156.4e6}{km^3}, antara 300 dan \qty{400}{m}: $300 + (156.4 - 117.9)/42.6 \times 100
= \qty{390}{m}$.
\textbf{23.} $38.7 + 0.90 \times 9.4 = \qty{47.2}{\percent}$ daratan masa kini.
\textbf{24.} Menggandakan laju pemekaran akan menaikkan laut sekitar \textbf{\qty{390}{m}} dan
menenggelamkan hampir separuh daratan masa kini.
\end{solution}
